\begin{abstract}
Conventional wisdom suggests quality-first curation dominates at small scale while diversity-first should take over at large scale as quality saturates---but through systematic GPT-2 training experiments across four log-spaced scales (10K--10M tokens), we find quality-first curation universally outperforms diversity-first by +1.5 percentage points (pp) to +5.0pp (Cohen's $d=0.76$--$2.48$, medium-to-very-large effect sizes), contradicting the predicted reversal. Data curation for foundation model training combines quality filtering and diversity sampling, yet optimal sequential ordering remains unclear. We quantify quality saturation (slope $-2.62$pp/log-scale, $p=0.008$) and model diversity persistence based on Feature Activation Coverage literature (simulated trajectory pending full validation, slope $+2.5$pp/log-scale). Our compositional ordering experiments show evidence for non-additive interactions that extend Data Mixing Laws by demonstrating order matters for sequential filtering steps. The underlying mechanism appears to be quality-gated diversity: diversity sampling shows effectiveness primarily on quality-filtered subsets, not raw noisy corpora. These findings suggest prescriptive guidance for practitioners working with 100M--1B parameter models on web corpora at 10K--10M scale---apply quality filtering (V-Information) before diversity sampling (Feature Activation Coverage)---and highlight the importance of compositional ordering in multi-stage data curation pipeline design.
\end{abstract}
